\documentclass[a4paper, 11pt]{article}

% -- Coments
\usepackage{verbatim}

% ----- Fonts -----
% -- Fuente --
% \usepackage{fontspec}
% \setmonofont{JetBrainsMono Nerd Font}  

% -- Color --
\usepackage{xcolor}
%\definecolor{azul}{RGB}{00,33,99}
\definecolor{azul}{RGB}{35,72,180}

% -- Page Margin --
\usepackage[margin=1in]{geometry}

% -- Espaciados --
\newcommand{\Pspace}{0.5cm}
\newcommand{\Aspace}{0.2cm}

% -- Columnas --
\usepackage{multicol}

% -- Imagenes --
\usepackage{graphicx}
\usepackage{float}

% -- Matemáticas --
\usepackage{amsmath, amssymb}

% -- Gráficas --
\usepackage{pgfplots}
\pgfplotsset{compat=1.18}


\title
{
    Contexto Regional 2026-2 \\
    Historia General de Sonora vol. 1 - María Elisa Villalpando \\

    Profesor José Moreno Vega
}

\begin{document}

\maketitle

\begin{center}
    \begin{tabular}{r|l}
        \textbf{Expediente} & \textbf{Nombre} \\ \hline
        219208106 & Bórquez Guerrero Angel Fernando \\
    \end{tabular}
\end{center}

{\Large 1. ¿Cuál es el argumento central del autor?} \\
Habla sobre cómo la cultura de los seris y su historia han sobrevivido en el estado a pesar de ser una cultura principalmente nómadas y no prácticar la agricultura, no sin antes recalcar su comprehensión por el arte de la ceramica, algo lo cual se encuentra por toda la expansión de su terrotorio, los cuales abarca desde la parte central del desierto de sonora hasta la costa del actual guaymas. Considerando todos estos puntos, creo que lo que quiere dar a entender el autor es su increible capacidad de resilenecía y resistencía ante los intentos de absorción de otro grupos índigenas o de las misiones.

{\Large 2. ¿Qué evidencias utiliza para respaldarlo?} \par
Vasijas, imágenes, y distintos objetos de caza, son los príncipales expositores de su cultura, así como, diferentes registros de misioneros como Eusebio Francisco Kino.

{\Large 3. ¿Cómo se relaciona  esta lectura con los temas vistos en las clases anteriores?} \par
Se podría decir que realmente lo único que permitia la expansión de culturas antiguas eran las condiciones climatológicas favorables del territorio de mesoamérica, ya que al intentar desplazarse a un territorio no parecido al suyo, siempre retrocedian ante las extremas o "desfavorables" condiciones de Aridoamérica y simplemente dejaban ser a las demás culturas.

\end{document}
